% New English translation from checked cote 25, site 70 (PDF page 71).
% AI-translated by GPT-6 Astra (OpenAI, ChatGPT Pro mode).
\scanpage{25/70}
\mspage{65}

is indeed independent of the basis\hand{/}chosen. In the case m=1, we have
the obvious intrinsic interpretation mentioned in 5.2. With this
linguistic convention, 5.3 remains valid as it stands, and likewise 5.4.
The first part of remark 5.5 admits a \var{\xout{generalization}}{}{typewriter strike-out}{ega25_25_070-v0001} \var{\inserted[typed]{\var{\xout{variant for the}}{}{typewriter strike-out}{ega25_25_070-v0003}}}{\var{\xout{variant for the}}{}{typewriter strike-out}{ega25_25_070-v0003}}{insertion}{ega25_25_070-v0002}
\var{\inserted[typed]{generalization to the}}{generalization to the}{insertion}{ega25_25_070-v0004}
\var{\xout{to the}}{}{typewriter strike-out}{ega25_25_070-v0005} case of arbitrary m: when $F_s$ is \var{\xout{\unsure{Sm}}}{}{typewriter strike-out}{ega25_25_070-v0006} \var{\inserted[typed]{$(S_m)$}}{$(S_m)$}{insertion}{ega25_25_070-v0007}, then the regularity condition
considered for $\hand{\phi}_{\hand{\xi}}$ is expressed purely in terms of dimension.
The second part of remark 5.5 holds as it stands for arbitrary
m. \unsure{Theorem} 5.6 extends as it stands, and likewise 5.7.

\uline{Proposition} 6.1: read dim $f(X_i) \hand{\geqslant}$ m+1, and
later dim $f(Z) \hand{\leqslant}$ m-1.

The general spiels of 7.1 hold as they stand for arbitrary
m. 7.2 and 7.3 mutatis mutandis (beware in 7.2
of the \unsure{notation} m, which is confusing here); on the other hand, in the proof
of 7.4 we no longer need to proceed step by step, but can
directly take a linear section of codimension m=n.

In 8.2, replace \unsure{the}\var{\xout{hypothesis}}{}{typewriter strike-out}{ega25_25_070-v0008} condition dim $Y_{\hand{\xi}}$ \hand{$>$} dim X-\unsure{1}
by dim $Y_{\hand{\xi}}$ \hand{$>$} dim X - \var{\over{1}{m}}{m}{writing over}{ega25_25_070-v0009}, and the hypothesis dim $f(X_i) \hand{>} 0$ by
dim $f(X_i) \hand{\geqslant}$ m. Make an analogous modification in the remainder of 8.2. Since
8.3 presents an example, there is no reason to change it, so we
shall keep m=1 there. I leave to you the exercise, doubtless requiring
\hand{/}a little more care, of finding the right statements for arbitrary m
corresponding to 8.4, 8.5 and 8.6 (pages \var{\inserted{30,}}{30,}{insertion}{ega25_25_070-v0010}31,32)\hand{;}
\hand{it is probably not indispensable, moreover, to}
\hand{do this exercise if doing it bores you.}
\tnote{25-070-1}{The internal m/p doubt in Théorème, the a doubt in notation, and the final-letter doubt in la are carried by the corresponding English words. The colloquial s'il te barbe is rendered if it bores you. The source's references to 8.4, 8.5, 8.6 and its page-number insertion are retained.}

\mspage{70}
